% Chapter XI, §106 (b)--(c)
% Chandrasekhar, "Hydrodynamic and Hydromagnetic Stability"
% Book pages 510--512
% Continues from §106 (a); equation numbering from (197)

%%% ------------------------------------------------------------------
%%% (b) Two uniform fluids separated by a horizontal boundary
%%%     (magnetic field parallel to the streaming)
%%% ------------------------------------------------------------------
\subsection*{(b) The case of two uniform fluids in relative horizontal
  motion separated by a horizontal boundary}

We shall apply the foregoing equations to the case of two uniform
fluids of densities $\rho_1$ and $\rho_2$ separated by a horizontal boundary
at $z=0$.  Let the velocities of streaming of the two fluids be $U_1$
and~$U_2$.

In the two regions of constant $\rho$ and~$U$, equations~\eqref{eq:11-190}
and~\eqref{eq:11-196} give
\begin{equation}
\label{eq:11-197}
\zeta = 0
\end{equation}
and
\begin{equation}
\label{eq:11-198}
\left[p(n+k_x\,U) - \frac{k_x^{2}}{n+k_x\,U}\,
\frac{\mu H^{2}}{4\pi}\right](D^{2}-k^{2})w = 0,
\end{equation}
where we have suppressed the subscripts distinguishing the two fluids.
The solutions in the two regions $z<0$ and $z>0$ can be written in the
forms (cf.\ equations~(24) and~(25))
\begin{equation}
\label{eq:11-199}
w_1 = A(n+k_x\,U_1)\,e^{+kz}\quad (z<0)
\end{equation}
and
\begin{equation}
\label{eq:11-200}
w_2 = A(n+k_x\,U_2)\,e^{-kz}\quad (z>0).
\end{equation}
By this choice of the solutions, the condition requiring the continuity
of $w/(n+k_x\,U)$, as well as the conditions at infinity, have been met.
It remains to satisfy the further condition which follows from
equation~\eqref{eq:11-196} by integrating it across the interface at
$z=0$.  This condition is, clearly,
\begin{equation}
\label{eq:11-201}
\Delta_0\!\bigl\{p(n+k_x\,U)Dw\bigr\}
= k_x^{2}\,\Delta_0\!\!\left(\frac{Dw}{n+k_x\,U}\right)
  + gk^{2}\Delta_0(p)\!\left(\frac{w}{n+k_x\,U}\right)_{\!0},
\end{equation}
where $\Delta_0(f)$ denotes the jump a quantity~$f$ experiences at $z=0$
and $[w/(n+k_x\,U)]_0$ is the unique value which this quantity has at
$z=0$.

Now applying the condition~\eqref{eq:11-201} to the solutions given by
equations~\eqref{eq:11-199} and~\eqref{eq:11-200}, we obtain the
characteristic equation
\begin{equation}
\label{eq:11-202}
\rho_1(n+k_x\,U_1)^{2}+\rho_2(n+k_x\,U_2)^{2}
= gk(\rho_1-\rho_2)+k_x^{2}\,\frac{\mu H^{2}}{2\pi}.
\end{equation}

On comparing this equation with equation~(26), we observe that in this
instance the effect of the magnetic field is equivalent to that of a
surface tension of amount
\begin{equation}
\label{eq:11-203}
T_{\text{eff}} = \frac{\mu H^{2}}{2\pi}\,\frac{k_x^{2}}{k^{2}}.
\end{equation}
This is exactly the same expression we found in the development of the
Rayleigh--Taylor instability in Chapter~X, \S\,97, equation~(238).

The roots of equation~\eqref{eq:11-202} are (cf.\ equation~(30))
\begin{equation}
\label{eq:11-204}
n = -k_x(\alpha_1\,U + \alpha_2\,U_2) \pm
\left[gk(\alpha_1-\alpha_2) + k_x^{2}\,\frac{\mu H^{2}}{2\pi}
  - k_x^{2}\,\alpha_1\,\alpha_2(U_1-U_2)^{2}\right]^{1/2}\!.
\end{equation}

From this equation it follows that \textit{a uniform magnetic field acting
parallel to the direction of streaming will suppress the
Kelvin--Helmholtz instability if}
\begin{equation}
\label{eq:11-205}
\alpha_1\,\alpha_2\,(U_1-U_2)^{2}
\leqslant \frac{\mu H^{2}}{2\pi(\rho_1+\rho_2)},
\end{equation}
\textit{i.e.\ if the relative speed does not exceed the root-mean-square
Alfv\'en speed in the two media.}

%%% ------------------------------------------------------------------
%%% (c) The effect of a magnetic field transverse to the direction
%%%     of streaming
%%% ------------------------------------------------------------------
\subsection*{(c) The effect of a magnetic field transverse to the
  direction of streaming}

When the impressed magnetic field is transverse to the direction
of~$U$, the relevant perturbation equations are
\begin{equation}
\label{eq:11-206}
\mathrm{i}p(n+k_x\,U)u + p(DU)w
+ \frac{\mu H}{4\pi}\,(\mathrm{i}k_x\,h_x - \mathrm{i}k_y\,h_y)
= -\mathrm{i}k_x\,\delta p,
\end{equation}
\begin{equation}
\label{eq:11-207}
\mathrm{i}p(n+k_x\,U)v = -\mathrm{i}k_y\,\delta p,
\end{equation}
\begin{equation}
\label{eq:11-208}
\mathrm{i}p(n+k_x\,U)w
- \frac{\mu H}{4\pi}\,(\mathrm{i}k_x\,h_x - Dh_x)
= -D\delta p - g\,\delta p,
\end{equation}
and
\begin{equation}
\label{eq:11-209}
\mathrm{i}(n+k_x\,U)\mathbf{h}
= \mathrm{i}k_x\,H\,\mathbf{u} + h_x\,DU.
\end{equation}
These equations replace equations~\eqref{eq:11-172}--\eqref{eq:11-175};
and the remaining equations are unaltered.

Treating the present set of equations in the same manner as in
\S\,(a) above, we now obtain in place of~\eqref{eq:11-196}
\begin{align}
\label{eq:11-210}
D\!\bigl\{p(n+k_x\,U)Dw - pk_x(DU)w\bigr\}
&= k^{2}p(n+k_x\,U)w
  + k_x^{2}\,\frac{\mu H^{2}}{4\pi}\,
    D\!\left\{\frac{1}{n+k_x\,U}
      \!\left(\mathrm{i}u + \frac{DU}{n+k_x\,U}\,w\right)\right\}
\notag\\
&\quad
  - \frac{k_x^{2}\,w}{n+k_x\,U}\,\frac{\mu H^{2}}{4\pi}\,w
  + gk^{2}\,\frac{Dp}{n+k_x\,U}\,w.
\end{align}

We observe that equation~\eqref{eq:11-210} reduces to hydrodynamic
equation~(16) of \S\,100 in case $k_x=0$.  Therefore, \textit{the
development of the Kelvin--Helmholtz instability, in the direction of
the streaming, is uninfluenced by the presence of the magnetic field in
the transverse direction.}
